\catcode`\{=1 \catcode`\}=2 \catcode`\$=3 \catcode`\&=4 \catcode`\#=6 \catcode`\^=7 \catcode`\_=8 \catcode`\ =10 \catcode`\%=14
% lockstep prelude: every case \input's this file first. It boots the raw
% INITEX catcode table, enables PDF mode, turns all behaviour tracing to
% maximum, and provides \lsshipbox for shipping a box register.
% Notes (verified against pdfTeX 1.40.29):
% - \tracingall does not exist in -ini mode (it is a plain.tex macro), so
%   the equivalent switches are set individually here.
% - Both engines run with -ini -etex ("entering extended mode" in the log),
%   so the full e-TeX tracing set below is active too.
% - \tracingstats stays 0: it reports engine-internal memory use, not
%   behaviour, so no second implementation could ever match it.
% - Box dumps come from \tracingoutput (the same show_box routine \showbox
%   uses). An explicit \showbox/\showthe ends the run under the mandated
%   -halt-on-error flag (the "! OK." pseudo-error is fatal), so cases must
%   use \message with \the/\meaning instead of any \show* command.
\pdfoutput=1
\tracingonline=1
\tracingcommands=3
\tracingstats=0
\tracingpages=1
\tracingoutput=1
\tracinglostchars=1
\tracingmacros=2
\tracingparagraphs=1
\tracingrestores=1
\tracingassigns=1
\tracinggroups=1
\tracingifs=1
\tracingscantokens=1
\tracingnesting=2
% plain \tracingall sets \errorcontextlines=\maxdimen; \maxdimen is
% undefined in -ini, so spell its value (2^30-1) explicitly.
\errorcontextlines=1073741823
\showboxdepth=10000
\showboxbreadth=10000
% Ship counter and per-shipout marker (kept for debugging; run.py splits
% boxes on "Completed box being shipped out", not on this marker). Cases
% must not emit that text and must not touch \count250 (\lsshipcount).
\countdef\lsshipcount=250
\let\lsprimshipout=\shipout
\def\lsshipbox#1{\global\advance\lsshipcount by1\message{LOCKSTEP-BOX \the\lsshipcount}\lsprimshipout\box#1}
